\textbf{Bright Siren Searches Enabled by Rubin Observatory:}
% Rubin ToO - drafted
I will lead the observing planning, execution, and analysis of the next bright standard siren accessible to \Rubin, along with other \ToO's pursued during the \LSST, continuing my role as the \Rubin\ \ToO\ observer throughout the transition to operations.
% This leadership is enabled by my specialized expertise in the alert processing of the \ToO\ system, my intimate knowledge of the technical strengths and limitations of the observatory, and my sustained development of the \ToO\ system within the wider \Rubin\ ecosystem.
While it is difficult to predict the timing of the next bright siren, I will continue to develop and deploy improvements to the \ToO\ system to improve response latency, observing strategy, and system automation to maximize the forthcoming opportunity to detect and characterize these optical counterparts to obtain a well-constrained measure of \Ho.
% Lead DESGW IRX analyses - observation planning, image processing, and transient analysis - drafted

In parallel to operating the \Rubin\ \ToO\ system, I will support the analysis of \DESGW\ candidates identified during \LVK\ O4, and lead the analysis of \DESGW\ candidates identified during \LVK\ intermediate runs. This work positions me to sustain and support \Rubin's \ToO\ program with additional imaging in other filters or to extended time-frames using a similarly capable photometric imager. This increased temporal sampling and color characterization will play an important role, especially as other priorities of the \LSST\ are balanced against the \ToO\ program.
% Other observing programs with spectroscopy to compliment the Rubin alert stream - in progress
As \LSST's \ToO\ program matures, I will work to enhance counterpart characterization by securing spectroscopic follow-up time on other instruments, moving beyond photometric classification to confirm spectral features that can inform the scope of \Rubin's \ToO s and constrain the physical properties of the underlying transients and their host environments. 
% This includes an ongoing program with \JGEM\ using \PFS, a new program using \LLAMAS\ to characterize promising candidates from the \LSST\ alert stream, and a future program using data from the BOOMBOX spectrograph from the Via project\TODO{Double check if applying to a Carnegie institution, otherwise swap with SOAR}.

% % Simulations of O5, 3G, and LISA source identification using Rubin
% Looking further ahead, I will begin preparing \Rubin's \ToO\ infrastructure for multi-messenger astrophysics of the future, including \threeG\ ground-based gravitational-wave detectors, the space-based \GW\ detector \LISA, new classes of high-energy neutrinos from IceCube gen-2, and the forthcoming long-baseline neutrino experiments, \DUNE\ and \hyperk. This era will usher in larger event rates, more precise localizations, and new source classes (including massive black hole mergers), that demand \ToO\ strategies tailored for a new status quo. This work will focus on identifying the optimal observing strategies for \Rubin\ discovery using \texttt{rubin\_sim} simulations along with simulated light curves, updating the alert-handling infrastructure for rapid processing, and strengthening follow-up partnerships across a variety of other observatories and experiments. Enhancing the detection of counterparts associated with \GW\s from \threeG\ and space-based detectors will extend the standard siren lever arm for precision cosmology, enabling measurements of $\Omega_m$, $w$, and other dark energy models.
